% Part: computability
% Chapter: computability-theory
% Section: coding-computations

\documentclass[../../../include/open-logic-section]{subfiles}

\begin{document}

\olfileid{cmp}{thy}{cod}
\olsection{ترميز الحسابات}

في كل نموذج من نماذج الحساب تتأتى الأمور الآتية:
\begin{enumerate}
\item ضبط \emph{تعريفات} الدوال القابلة للحساب في وصف منتظم؛
  كمواصفات آلات تورنغ، أو التعريفات العودية، أو البرامج المكتوبة
  بإحدى لغات البرمجة.
\item تدوين حساب الدالة، متى عُرف تعريفها ودخلها، في سجل تام. فحساب
  آلة تورنغ، مثلًا، تسجله متتالية تهيئاتها: حالة الآلة وما على
  الشريط في كل خطوة.
\item فحص سجل يُدّعى أنه سجل حساب، للتحقق من مطابقته لحساب
  دالة قابلة للحساب ذات تعريف معلوم عند دخل معين، على أن تكون الدالة معرّفة
  عند هذا الدخل، أي أن ينتهي الحساب.
\item استخراج قيمة الدالة عند ذلك الدخل من سجل
  الحساب التام؛ ففي آلة تورنغ المتوقفة، مثلًا، تؤخذ القيمة من محتوى الشريط
  عند الخطوة الأخيرة.
\end{enumerate}

وبالترميز نجعل لكل وصف لدالة قابلة للحساب \emph{فهرسًا} عدديًا،
بحيث تُسترد تعليماته من فهرسه بطريقة قابلة
للحساب. وكذلك يُجعل سجل الحساب الكامل عددًا واحدًا. ومن هذا
تنشأ العلاقة الحسابية «$s$ يرمز سجل حساب الدالة ذات الفهرس~$e$
عند الدخل~$x$»، والدالة «خرج متتالية الحساب التي رمزها~$s$». وهما قابلتان
للحساب، بل عوديتان بدائيتان.

وعلى هذه الحقيقة الأساسية تبتني نتائج مهمة بعيدة الأثر
في قابلية الحساب؛ فنثبتها من غير أن يختص البرهان بنموذج الحساب
الذي اخترناه.

\end{document}
